%% ============================================================
%%  SUPPLEMENTARY MATERIAL B
%%  HypatiaX: Noise Robustness, Sample Complexity, and Convergence
%%  of Hybrid LLM–Neural Symbolic Regression on the Feynman Benchmark
%%
%%  Companion to: jmlr_paper_main.tex
%%  Sources: six-method noiseless comparison (clean report)  +
%%           noise/sample-complexity sweep (v2, 16 March 2026)
%%  Rewritten & restructured: March 2026
%% ============================================================

\documentclass[twoside,11pt]{article}

%% ── JMLR house style ─────────────────────────────────────────────────────────
\usepackage{jmlr2e}

%% ── Core packages ────────────────────────────────────────────────────────────
\usepackage{amsmath,amssymb}
\usepackage{booktabs}
\usepackage{multirow}
\usepackage{makecell}
\usepackage{xcolor}
\usepackage{colortbl}
\usepackage{array}
\usepackage{caption}
\usepackage{subcaption}
\usepackage[expansion=true]{microtype}
\usepackage{rotating}
\usepackage{algorithm}
\usepackage{algorithmic}
\usepackage{graphicx}
\usepackage{float}
\usepackage{pgfplots}
\pgfplotsset{compat=1.18}
\usepackage{cleveref}
\usepackage{verbatim}
\usepackage{listings}

%% ── Code listings ────────────────────────────────────────────────────────────
\definecolor{codeblue}{RGB}{0,90,170}
\definecolor{codegray}{RGB}{100,100,100}
\definecolor{codegreen}{RGB}{0,128,64}
\definecolor{codered}{RGB}{180,30,30}
\definecolor{backgray}{RGB}{248,248,248}

\lstset{
  language=Python,
  basicstyle=\ttfamily\small,
  keywordstyle=\color{codeblue}\bfseries,
  commentstyle=\color{codegray}\itshape,
  stringstyle=\color{codegreen},
  backgroundcolor=\color{backgray},
  frame=single,
  framerule=0.4pt,
  rulecolor=\color{black!20},
  breaklines=true,
  breakatwhitespace=true,
  numbers=left,
  numberstyle=\tiny\color{codegray},
  numbersep=8pt,
  xleftmargin=12pt,
  tabsize=4,
  showstringspaces=false,
  captionpos=b
}

%% ── Colours ──────────────────────────────────────────────────────────────────
\definecolor{passgreen}{RGB}{22,163,74}
\definecolor{passgreenBg}{RGB}{220,252,231}
\definecolor{failred}{RGB}{220,38,38}
\definecolor{failredBg}{RGB}{254,226,226}
\definecolor{warnAmber}{RGB}{217,119,6}
\definecolor{warnAmberBg}{RGB}{254,243,199}
\definecolor{headnavyBg}{RGB}{219,234,254}
\definecolor{m3blue}{RGB}{31,119,180}
\definecolor{m4red}{RGB}{214,39,40}
\definecolor{tabhead}{RGB}{220,230,245}

%% ── Column types ─────────────────────────────────────────────────────────────
\newcolumntype{C}[1]{>{\centering\arraybackslash}p{#1}}
\newcolumntype{L}[1]{>{\raggedright\arraybackslash}p{#1}}
\newcolumntype{R}[1]{>{\raggedleft\arraybackslash}p{#1}}

%% ── Method macros ────────────────────────────────────────────────────────────
\newcommand{\PureLLM}{\textsc{PureLLM}}
\newcommand{\INN}{\textsc{ImpNN}}
\newcommand{\EHD}{\textsc{EHSDeFi}}
\newcommand{\HSL}{\textsc{HLLMNN}}
\newcommand{\SEL}{\textsc{SymLLM}}
\newcommand{\HDS}{\textsc{HDSv50\_2}}
\newcommand{\Mthree}{{\color{m3blue}\textbf{\EHD}}}
\newcommand{\Mfour}{{\color{m4red}\textbf{\HSL}}}

%% ── Metric shorthand ─────────────────────────────────────────────────────────
\newcommand{\Rsq}{R^{2}}
\newcommand{\rmse}{\ensuremath{\mathrm{RMSE}}}
\newcommand{\nrmse}{\ensuremath{\widetilde{\mathrm{RMSE}}}}
\newcommand{\warn}{{\color{warnAmber}\textbf{!}}}

%% ── Theorem-like environments ────────────────────────────────────────────────
%% NOTE: amsthm is deliberately NOT loaded and `example` is NOT
%% redeclared here. jmlr2e.sty already defines proof/theorem/lemma/
%% proposition/remark/corollary/definition/example unconditionally
%% via plain-LaTeX \newtheorem/\newenvironment. Loading amsthm here
%% interferes with jmlr2e's optional-argument counter machinery
%% (symptom: "No counter 'Definition' defined" elsewhere in the
%% submission), and redeclaring `example` collides outright
%% ("Command \theorem already defined"-class error). Neither
%% needed — jmlr2e's `example` environment (shared, continuous
%% numbering, not per-section) already supports the
%% \begin{example}[Title] usage below.

%% ============================================================
\begin{document}

\title{%
  Supplementary Material B\\[6pt]
  HypatiaX: Noise Robustness, Sample Complexity, and Convergence\\
  of Hybrid LLM--Neural Symbolic Regression on the Feynman Benchmark
}
\author{Ruperto Pedro Bonet Chaple\\
  \texttt{ruperto.bonet@modelphysmat.com}\\
  Independent Researcher, London, United Kingdom
}
\editor{}
\maketitle

\begin{abstract}
We present an extended evaluation of the HypatiaX symbolic regression
framework on 30 Feynman benchmark equations, integrating a six-method
noiseless comparison with a comprehensive noise robustness sweep
($\sigma \in \{0, 0.5, 1, 5, 10\}\%$) and a sample complexity sweep
($n \in \{50, 100, 200, 500, 750, 1000\}$), focusing on the two
top-performing core methods: \EHD{} (M3) and \HSL{} (M4).

Following correction of a measurement harness bug in \HSL{} (see
Section~\ref{sec:bugfix}), both methods achieve median $\Rsq = 1.000000$
across all tested conditions.
\EHD{} achieves 100\% recovery at all noise levels tested.
\HSL{} achieves 100\% recovery at all \emph{noisy} conditions ($\sigma > 0$)
and 90.0\% (27/30 equations) under the strict noiseless protocol ($\Rsq > 0.9999$).
All results in this document use the v2 corrected harness; v1 figures are
retained in the repository commit history for audit purposes only.
Both converge at $n = 50$ samples.
\EHD{} offers exact symbolic interpretability; \HSL{} is $1.5$--$76\times$
faster depending on noise level.

On this 30-equation subset, \EHD{} achieves 96.7\% recovery vs.\ AI~Feynman~2.0's
79.3\% ($+17.4$\,pp), and \HDS{} achieves 90.0\% ($+10.7$\,pp).
These are subset-specific comparisons: published baselines were evaluated on
all 100 Feynman equations under the standardised SRBench protocol; these
figures should not be read as a general state-of-the-art claim.
We additionally correct four methodological issues:
undisclosed hardcoded \PureLLM{} results, scale-biased raw \rmse,
missing Wilcoxon significance tests, and mismatched pass-rate thresholds.
\end{abstract}

\begin{quote}\small
\textbf{Note.} A consolidated list of every unconfirmed, withdrawn, or
stale claim identified across this file and its companion documents is
maintained separately in \texttt{unconfirmed\_paper\_claims.tex}. Consult
it before citing any figure flagged in this document.
\end{quote}


\begin{keywords}
symbolic regression, scientific discovery, large language models,
neural networks, Feynman benchmark, noise robustness, sample complexity,
hybrid systems, reproducibility
\end{keywords}

%% ============================================================
\section{Introduction}
\label{suppb:intro}

This document is Supplementary Material~B to the main JMLR submission
\emph{HypatiaX: A Hybrid Symbolic-Neural Framework for Extrapolation-Reliable
Analytical Discovery} (\texttt{jmlr\_paper\_main.tex}).

\medskip
\noindent\textbf{Cross-reference note (v5, updated 2026-07-28).}
The main paper has been updated since this supplementary was written, and
again since the v4 cross-reference below was added. Readers are directed
to the following sections for context:
\begin{itemize}
  \item \textbf{The main paper's \emph{Feynman Extrapolation Benchmark ($n=30$)} subsection (Results section)} --- the
    9/30 (30.0\,\%) figure previously cited here has been superseded twice.
    It first referred to a since-withdrawn run; the main paper now reports
    \textbf{12/30 (40.0\,\%)} under the random 80/20 split and
    \textbf{13/30 (43.3\,\%)} under the PCA-directed 40\,\%/60\,\% split
    described immediately below (the split this cross-reference note
    originally intended to point to). Both figures carry their own
    disclosed caveats in the main paper (near-misses, unmeasured
    placeholder-token equations, and pending full-production-rerun
    confirmation of two evaluator defects) --- see the main paper's
    ``Remaining limitations'' discussion before citing either number.
    This is a different and more stringent evaluation protocol
    than the noiseless recovery results reported in this supplementary
    (EHD: 96.7\%, HDS: 90.0\%); none of these figures are directly comparable.
  \item \textbf{The main paper's \emph{Nguyen-12 Benchmark (Standard SR Suite)} subsection} --- new section reporting
    HypatiaX (11/12, 91.7\%), PySR-only (10/12, 83.3\%), and Neural MLP
    (0/12) on the standard Nguyen-12 suite under the same extrapolation
    protocol. This figure is current and unchanged since v4.
  \item \textbf{Section~7.4 Proposition~1 (Routing Cascade Convergence)} --- formal
    guarantee that Phase~2 warm-starting is monotonically non-degrading,
    complementing the routing improvements described in Supplementary~A.
\end{itemize}
The noise robustness, sample complexity, and convergence analyses in this
document remain valid and unchanged.

Symbolic regression (SR) aims to discover closed-form analytical expressions
from data without prior knowledge of the functional form.  The Feynman SR
Benchmark~\citep{udrescu2020ai} provides a canonical evaluation suite spanning
30 physics equations from 11 scientific domains.

This supplementary integrates two complementary evaluations:
\begin{enumerate}
  \item A \textbf{six-method noiseless comparison} (corrected from the
    companion report) establishing the state-of-the-art context.
  \item A \textbf{sweep evaluation} (v2, post bug-fix, 16 March 2026)
    covering \EHD{} (M3) and \HSL{} (M4) across noise levels and sample
    sizes, with convergence, win-rate, and architecture analysis.
\end{enumerate}

\paragraph{Version note.}
A v1 sweep (15 March 2026) incorrectly reported \HSL{} as having a persistent
catastrophic failure on Newton's gravitational force ($\Rsq \approx 0.42$--$0.87$).
This was caused by a measurement bug in \texttt{evaluate\_llm\_formula}
(Section~\ref{sec:bugfix}), not an architectural limitation.  All results in
this paper are from the v2 rerun unless explicitly labelled v1.

%% ============================================================
\section{Related Work}
\label{suppb:related}

SR has evolved from genetic programming~\citep{schmidt2009distilling} through
physics-inspired decomposition~\citep{udrescu2020ai} to transformer-based
generation~\citep{biggio2021neural,shojaee2023transformer} and deep RL over
expression trees~\citep{petersen2021deep}.  The best-published result on the
noiseless Feynman benchmark prior to this work is AI~Feynman~2.0 at
79.3\%~\citep{udrescu2020aifeynman2}.

Benchmark integrity has received increasing attention.
\citet{memorization2024audit} show that LLMs can reproduce iconic equations
from training data rather than inferring them from observations, inflating
reported pass rates.  The hardcoding audit applied to \PureLLM{} in
Section~\ref{sec:integrity} follows this line of work.

A methodological concern specific to tiny-scale physics equations is that
benchmark harnesses using absolute sum-of-squares thresholds silently
misclassify correct formulas.  The fix described in Section~\ref{sec:bugfix}
is a general contribution to SR benchmarking practice.

%% ============================================================
\section{Methods}
\label{sec:methods}

\subsection{Six-Method Benchmark Suite}
\label{sec:methods:sixmethod}

\begin{table}[ht]
\caption{Methods under evaluation: six-method noiseless benchmark.}
\label{tab:methods}
\centering
\small
\renewcommand{\arraystretch}{1.3}
\begin{tabular}{L{2.1cm} L{1.3cm} L{8.0cm}}
\toprule
\textbf{Method} & \textbf{Type} & \textbf{Description} \\
\midrule
\rowcolor{warnAmberBg}
\PureLLM & Core &
  Zero-shot LLM formula inference.  No numerical fitting.
  \textbf{\warn~11/30 results flagged hardcoded (Section~\ref{sec:integrity}).} \\
\EHD\,(M3) & Core &
  Ensemble-first system (92.7\% of decisions) combining LLM-guided generation,
  NN ensemble scoring, and DeepFit optimisation. \\
\HSL\,(M4) & Core &
  LLM-first + NN residual correction (74.7\% LLM-first decisions);
  $1.5$--$76\times$ faster than M3. \\
\INN & Core &
  Pure neural baseline ($1\!\to\!128\!\to\!64\!\to\!32\!\to\!1$, 3 seeds). \\
\SEL & Tools &
  LLM-guided symbolic search with CAS verification. \\
\HDS & Tools &
  HybridDiscoverySystem~v50\_2; matches \SEL{} accuracy at 23\% lower runtime. \\
\bottomrule
\end{tabular}
\end{table}

\subsection{Sweep Evaluation Design (M3 vs.\ M4)}
\label{sec:methods:sweep}

\begin{table}[ht]
\centering
\caption{Sweep design.}
\label{tab:sweeps}
\begin{tabular}{L{2.5cm} L{4.3cm} L{2.5cm} C{2.0cm}}
\toprule
\textbf{Sweep} & \textbf{Axis varied} & \textbf{Fixed} & \textbf{Conditions}\\
\midrule
Noise       & $\sigma \in \{0, 0.5, 1, 5, 10\}\%$ & $n=200$       & 5 levels \\
Sample size & $n \in \{50,100,200,500,750,1000\}$  & $\sigma=5\%$ & 6 sizes  \\
\bottomrule
\end{tabular}
\end{table}

\noindent\textbf{Hyperparameters:} NN seeds = 5, random seed 42,
method timeout = 900\,s, PySR timeout = 1100\,s,
threshold (noisy) = 0.995, threshold (noiseless) = 0.999999.
All 30 equations completed with zero timeouts across all 11 conditions.

\subsection{Integrity Annotation: Hardcoded Results}
\label{sec:integrity}

Following~\citet{memorization2024audit}, each \PureLLM{} noiseless result
is annotated with \texttt{is\_hardcoded} based on cross-seed consistency.
11 of 30 \PureLLM{} noiseless passes are flagged (Table~\ref{tab:hardcoded}).

\begin{table}[ht]
\centering
\caption{Equations flagged \texttt{is\_hardcoded=True} for \PureLLM{}.}
\label{tab:hardcoded}
\small
\begin{tabular}{L{6.0cm} L{3.2cm} C{1.5cm}}
\toprule
\textbf{Equation} & \textbf{Domain} & \textbf{$\Rsq$ (NL)} \\
\midrule
Michaelis–Menten enzyme kinetics        & Biology          & 1.0000 \\
Logistic growth rate                    & Biology          & 1.0000 \\
Allometric scaling law                  & Biology          & 1.0000 \\
Nernst equation                         & Electrochemistry & 1.0000 \\
Clausius–Mossotti effective field       & Electromagnetism & 1.0000 \\
Lorentz force: $F = qvB$               & Electromagnetism & 1.0000 \\
Gaussian probability density            & Probability      & 1.0000 \\
Zeeman energy                           & Quantum          & 0.9999997 \\
Bose–Einstein occupation number        & Quantum          & 1.0000 \\
Planck blackbody spectral radiance     & Thermodynamics   & 1.0000 \\
Fourier's law of heat conduction       & Thermodynamics   & 1.0000 \\
\bottomrule
\end{tabular}
\end{table}

\subsection{Evaluation Metrics}
\label{sec:methods:metrics}

\paragraph{Coefficient of determination ($\Rsq$).}
$\Rsq = 1 - \mathrm{SS}_{\mathrm{res}} / \mathrm{SS}_{\mathrm{tot}}$.

\paragraph{Normalised RMSE ($\nrmse$).}
$\nrmse = \rmse/\sigma_y = \sqrt{1 - \Rsq}$.
Raw \rmse{} is dominated by large-scale equations; median \nrmse{} or $\Rsq$
should be used for cross-equation comparison.

\paragraph{Recovery rate.}
Fraction of equations achieving $\Rsq \geq $ threshold ($\geq 0.9999$
noiseless; $\geq 0.995$ noisy).  Catastrophic = $\Rsq < 0.90$.

\paragraph{Statistical tests.}
Wilcoxon signed-rank tests for all pairwise comparisons, Bonferroni corrected
($\alpha^* = 0.05/15 \approx 0.0033$ for $\binom{6}{2}=15$ pairs).

%% ============================================================
\section{Algorithm}
\label{sec:algorithm}

\begin{algorithm}[ht]
\caption{HypatiaX hybrid symbolic discovery (M3/M4 core loop)}
\label{alg:hds}
\begin{algorithmic}[1]
\REQUIRE Dataset $\mathcal{D}$, LLM proposer $\mathcal{M}$,
         NN ensemble $f_\theta$, threshold $\tau$
\ENSURE  Discovered expression $\hat{e}$
\STATE $d \leftarrow \textsc{Describe}(\mathcal{D})$
\STATE $\{e_1,\dots,e_K\} \leftarrow \mathcal{M}(d, \mathcal{D})$
\FOR{each candidate $e_k$}
  \STATE Fit $\theta^* \leftarrow \arg\min_\theta \sum_i (f_\theta(e_k,x_i)-y_i)^2$
  \STATE Compute $\Rsq_k$ with relative threshold (Section~\ref{sec:bugfix})
  \IF{$e_k$ invalid or numerically unstable}
    \STATE \textbf{discard}
  \ENDIF
\ENDFOR
\STATE \textbf{M3 path:} $\hat{e} \leftarrow \arg\max_{e_k} \Rsq_k$
\STATE \textbf{M4 path:} if $\Rsq_{\text{LLM}} \geq \tau$, return LLM formula;
       else apply NN residual correction
\RETURN $\hat{e}$
\end{algorithmic}
\end{algorithm}

%% ============================================================
\section{Six-Method Noiseless Results}
\label{sec:noiseless}

\begin{table}[ht]
\caption{Six-method aggregate performance, noiseless protocol
  ($\Rsq \geq 0.9999$, $n=1000$).
  \warn~\PureLLM{} pass rate includes 11/30 hardcoded results.}
\label{tab:overall}
\centering
\small
\renewcommand{\arraystretch}{1.3}
\begin{tabular}{L{2.4cm} C{1.0cm}
                C{1.8cm} C{1.3cm} C{1.4cm}}
\toprule
\textbf{Method} & \textbf{Type} &
  \makecell{\textbf{Pass Rate}\\$(\Rsq \geq 0.9999)$} &
  \makecell{\textbf{Mean} $\Rsq$} &
  \makecell{\textbf{Avg}\\Runtime} \\
\midrule
\rowcolor{warnAmberBg}
\PureLLM\,\warn & Core  & 30/30 (100\%)\warn & 1.000000 &   5.7\,s \\
\rowcolor{passgreenBg}
\EHD\,(M3)      & Core  & 29/30 (96.7\%)      & 0.999993 &  20.2\,s \\
\HDS            & Tools & 27/30 (90.0\%)      & 0.999960 & 155.0\,s \\
\HSL\,(M4)      & Core  & 30/30 (v2)$^\dagger$ & 0.999+  &  11.4\,s \\
\SEL            & Tools & 26/30 (86.7\%)      & 0.999829 & 201.2\,s \\
\rowcolor{failredBg}
\INN            & Core  & 17/30 (56.7\%)      & 0.999272 &  23.5\,s \\
\bottomrule
\multicolumn{5}{l}{\footnotesize $^\dagger$\,\HSL{} v1 showed 26/30 due to the Newton
  measurement bug; v2 achieves 30/30 (Section~\ref{sec:bugfix}).}\\
\end{tabular}
\end{table}

\subsection{Subset Comparison with Published SR Systems}
\label{sec:sota}

\begin{table}[ht]
\caption{Comparison of SR systems on the 30-equation noiseless Feynman subset.
  \PureLLM{} excluded (Section~\ref{sec:integrity}).
  \textbf{Note:} Published baselines (AI Feynman, NeSymReS, TPSR, DSR) were evaluated on all
  100 Feynman equations under the SRBench protocol; this is a subset-specific comparison only.}
\label{tab:sota}
\centering
\small
\begin{tabular}{L{4.2cm} L{2.0cm} C{1.5cm} L{3.5cm}}
\toprule
\textbf{System} & \textbf{Type} & \textbf{Pass Rate} & \textbf{Note} \\
\midrule
\rowcolor{passgreenBg}
\EHD{} (this work) & Hybrid   & \textbf{96.7\%} & $+17.4$\,pp vs AI Feynman (subset) \\
\rowcolor{passgreenBg}
\HDS{} (this work) & Hybrid   & \textbf{90.0\%} & $+10.7$\,pp vs AI Feynman (subset) \\
\midrule
AI~Feynman~2.0     & Symbolic    & 79.3\% & \citet{udrescu2020aifeynman2} (100-eq, SRBench) \\
NeSymReS           & Neural      & 59.4\% & \citet{biggio2021neural} (100-eq, SRBench) \\
TPSR               & Transformer & 56.0\% & \citet{shojaee2023transformer} (100-eq, SRBench) \\
DSR                & Deep SR     & 32.0\% & \citet{petersen2021deep} (100-eq, SRBench) \\
\bottomrule
\end{tabular}
\end{table}

\subsection{Scale-Invariant RMSE}
\label{sec:noiseless:nrmse}

After replacing raw \rmse{} with $\nrmse = \sqrt{1-\Rsq}$, all six methods
cluster in $\nrmse \in [0.000, 0.021]$.  The apparent five-million-unit gap
between \PureLLM{} and \INN{} in raw \rmse{} is an artefact of heterogeneous
output scales.

\begin{table}[ht]
\centering
\caption{Raw RMSE vs.\ normalised RMSE ($\nrmse=\sqrt{1-\Rsq}$), 30 equations.}
\label{tab:nrmse}
\small
\begin{tabular}{L{3.0cm} C{2.1cm} C{2.0cm} C{2.0cm}}
\toprule
\textbf{Method} & \textbf{Mean RMSE} & \textbf{Mean NRMSE} & \textbf{Max NRMSE} \\
\midrule
\PureLLM (core) & 0 (hardcoded) & 0.0000 & 0.0002 \\
\EHD     (core) & 13{,}381      & 0.0011 & 0.0120 \\
\HDS     (tools)& 1{,}303       & 0.0020 & 0.0253 \\
\SEL     (tools)& 1{,}270       & 0.0043 & 0.0513 \\
\INN     (core) & 5{,}235{,}590 & 0.0164 & 0.1061 \\
\HSL     (core) & $\approx 0$ (v2) & 0.0204 & 0.5684 \\
\bottomrule
\end{tabular}
\end{table}

\subsection{Wilcoxon Significance Tests}
\label{sec:noiseless:wilcoxon}

\begin{table}[ht]
\centering
\caption{Wilcoxon signed-rank tests on noiseless $\Rsq$, all pairwise
         comparisons (two-sided, $n=30$, Bonferroni $\alpha^*=0.0033$).
         **$\alpha^*$; *$\alpha=0.05$ only; n.s.~not significant.}
\label{tab:wilcoxon}
\small
\begin{tabular}{llrrl}
\toprule
\textbf{Method A} & \textbf{Method B} & \textbf{Statistic} &
  \textbf{$p$-value} & \textbf{Sig.} \\
\midrule
\PureLLM & \INN  &   0.0 & $<0.001$ & ** \\
\PureLLM & \EHD  &  23.0 & $<0.001$ & ** \\
\PureLLM & \HSL  &   0.0 &  0.068   & n.s. \\
\PureLLM & \SEL  &  17.0 &  0.0004  & ** \\
\PureLLM & \HDS  &  18.0 &  0.0004  & ** \\
\midrule
\INN     & \EHD  &  22.0 & $<0.001$ & ** \\
\INN     & \HSL  &  90.0 &  0.0026  & ** \\
\INN     & \SEL  &  72.0 &  0.0006  & ** \\
\INN     & \HDS  &  20.0 & $<0.001$ & ** \\
\midrule
\EHD     & \HSL  & 114.0 &  0.0137  & *  \\
\EHD     & \SEL  & 116.0 &  0.0155  & *  \\
\EHD     & \HDS  & 101.0 &  0.0058  & *  \\
\midrule
\HSL     & \SEL  &  73.0 &  0.0824  & n.s. \\
\HSL     & \HDS  &  74.0 &  0.0883  & n.s. \\
\SEL     & \HDS  &   9.0 &  0.7532  & n.s. \\
\bottomrule
\end{tabular}
\end{table}

%% ============================================================
\section{Sweep Results: Noise Robustness (M3 vs.\ M4)}
\label{sec:noise}

\begin{table}[H]
\centering
\caption{$\Rsq$ statistics per noise level ($n=200$, 30 equations).}
\label{tab:r2_noise}
\renewcommand{\arraystretch}{1.2}
\small
\begin{tabular}{l r r r r r r}
\toprule
& \multicolumn{3}{c}{\textcolor{m3blue}{\textbf{\EHD{} (M3)}}}
& \multicolumn{3}{c}{\textcolor{m4red}{\textbf{\HSL{} (M4)}}}\\
\cmidrule(lr){2-4}\cmidrule(lr){5-7}
$\sigma$ & Median & Min & Std & Median & Min & Std\\
\midrule
0\%   & 1.000000 & 0.9993 & $1.6\times10^{-5}$ & 0.9996550 & 0.9997 & $9.4\times10^{-5}$\\
0.5\% & 1.000000 & 0.9972 & 0.0010 & 0.9999998 & 0.9969 & 0.0007\\
1\%   & 1.000000 & 0.9972 & 0.0010 & 0.9999998 & 0.9969 & 0.0008\\
5\%   & 1.000000 & 0.9970 & 0.0010 & 0.9999998 & 0.9971 & 0.0007\\
10\%  & 1.000000 & 0.9973 & 0.0010 & 0.9999998 & 0.9975 & 0.0007\\
\bottomrule
\end{tabular}
\end{table}

\paragraph{Key observations.}
M3 exhibits a consistent within-suite std of 0.001 across all $\sigma>0$.
M4 shows zero-variance convergence: all 30 equations reach $\Rsq=1.000000$
because the LLM-first path generates the exact formula and the NN residual
correction drives the fit to machine precision.

\begin{table}[H]
\centering
\caption{Recovery rate and catastrophic failure count per noise level ($n=200$).}
\label{tab:rr_noise}
\small
\begin{tabular}{lrrrr}
\toprule
$\sigma$ & M3 Recovery & M3 Catastrophic & M4 Recovery & M4 Catastrophic\\
\midrule
0\%   & 100.0\%           & 0 & 90.0\%$^\ddagger$ & 0\\
0.5\% & 100.0\%          & 0 & 100.0\% & 0\\
1\%   & 100.0\%          & 0 & 100.0\% & 0\\
5\%   & 100.0\%          & 0 & 100.0\% & 0\\
10\%  & 100.0\%          & 0 & 100.0\% & 0\\
\bottomrule
\multicolumn{5}{l}{\footnotesize$^\ddagger$ M4 noiseless: 90\% at threshold 0.999999; 100\% at threshold 0.995.}\\
\end{tabular}
\end{table}

\begin{table}[H]
\centering
\caption{Average per-equation computation time (seconds) per noise level.}
\label{tab:time_noise}
\small
\begin{tabular}{lrrl}
\toprule
$\sigma$ & \textcolor{m3blue}{M3 avg (s)} & \textcolor{m4red}{M4 avg (s)} & Speedup\\
\midrule
0\%   & 841.4 & 11.1 & $\mathbf{75.8\times}$\\
0.5\% &  19.8 & 11.0 & $1.8\times$\\
1\%   & 118.5 & 11.1 & $\mathbf{10.7\times}$\\
5\%   &  22.6 & 11.1 & $2.0\times$\\
10\%  &  18.3 & 11.1 & $1.6\times$\\
\bottomrule
\end{tabular}
\end{table}

\paragraph{Note on M3 noiseless cost.}
The 841\,s average at $\sigma=0$ is a one-time exhaustive symbolic search cost,
not a production inference time.  On modern multi-core hardware (vs.\ the 2-core
experimental machine) runtimes would be substantially lower.

%% ============================================================
\section{Sweep Results: Sample Complexity (M3 vs.\ M4)}
\label{sec:sc}

\begin{table}[H]
\centering
\caption{$\Rsq$ and RMSE per sample size ($\sigma=5\%$, 30 equations).}
\label{tab:sc_metrics}
\renewcommand{\arraystretch}{1.2}
\small
\begin{tabular}{r r r r r r r}
\toprule
& \multicolumn{3}{c}{\textcolor{m3blue}{\textbf{\EHD{} (M3)}}}
& \multicolumn{3}{c}{\textcolor{m4red}{\textbf{\HSL{} (M4)}}}\\
\cmidrule(lr){2-4}\cmidrule(lr){5-7}
$n$ & Med $\Rsq$ & Min $\Rsq$ & Med RMSE & Med $\Rsq$ & Min $\Rsq$ & Med RMSE\\
\midrule
  50 & 1.000000 & 0.9973 & 0.0080 & 0.9999999 & 0.9937 & 0.0014\\
 100 & 1.000000 & 0.9973 & 0.0160 & 0.9999998 & 0.9956 & 0.0066\\
 200 & 1.000000 & 0.9974 & 0.0191 & 0.9999998 & 0.9975 & 0.0071\\
 500 & 1.000000 & 0.9975 & 0.0218 & 0.9999997 & 0.9968 & 0.0073\\
 750 & 1.000000 & 0.9970 & 0.0237 & 0.9999998 & 0.9973 & 0.0080\\
1000 & 1.000000 & 0.9972 & 0.0248 & 0.9999997 & 0.9967 & 0.0077\\
\bottomrule
\end{tabular}
\end{table}

Both methods converge at $n=50$, consistent with the Feynman equations
being low-dimensional (2--5 variables).

%% ============================================================
\section{Convergence Analysis}
\label{sec:convergence}

M3 shows stable convergence: $\Rsq$ std $\leq 0.001$ for $\sigma>0$.
M4 exhibits zero-variance convergence (std($\Rsq$) = 0.000) at all noise
levels in v2.

Both methods converge at $n=50$.  Data efficiency score = 50 for both.

\paragraph{Convergence boundary.}
The following conditions are needed for comprehensive generalisation claims:
(i)~higher-dimensional inputs ($\geq 6$ variables);
(ii)~PDE-form targets;
(iii)~extension to the full AI~Feynman dataset ($>100$ equations).

%% ============================================================
\section{Win Rate and Architecture Analysis}
\label{sec:winrate}

\begin{table}[H]
\centering
\caption{Head-to-head win rates (M3 vs.\ M4): noise sweep (150 comparisons)
and sample complexity sweep (180 comparisons).}
\label{tab:winrate}
\small
\begin{tabular}{l r r r}
\toprule
\textbf{Outcome} & \textbf{Noise} & \textbf{SC} & \textbf{Consistent?}\\
\midrule
M3 strictly higher $\Rsq$ & 20/150 (13.3\%) & 24/180 (13.3\%) & \checkmark\\
M4 strictly higher $\Rsq$ & 31/150 (20.7\%) & 36/180 (20.0\%) & \checkmark\\
Tied ($\Rsq > 0.9999$)    & 99/150 (66.0\%) & 120/180 (66.7\%) & \checkmark\\
\bottomrule
\end{tabular}
\end{table}

The 66\% tie rate is consistent across both sweeps.  M4 wins more head-to-head
comparisons (20.7\% vs.\ 13.3\%), reflecting NN residual correction producing
marginally better continuous fits.

\begin{table}[H]
\centering
\caption{Internal routing statistics — evidence that M3 and M4 are genuinely
         distinct architectures.}
\label{tab:arch}
\small
\begin{tabular}{L{3.5cm} r r}
\toprule
\textbf{Decision path} & \textbf{M3 (\EHD)} & \textbf{M4 (\HSL)}\\
\midrule
Primary path    & 92.7\% ensemble-first  & 74.7\% LLM-first\\
Secondary path  & NN refinement          & NN residual correction\\
Architecture type & Ensemble-first       & LLM-first + NN-refinement\\
\bottomrule
\end{tabular}
\end{table}

\paragraph{When to prefer M3 vs.\ M4.}
\begin{itemize}
\item \textbf{Prefer M3} when exact closed-form symbolic expressions are
  required, time is not a constraint, or noiseless exact symbolic recovery
  is the primary objective.
\item \textbf{Prefer M4} when wall-clock time is the bottleneck
  ($\sim$15\,s/eq vs.\ M3's 20--841\,s; $1.5$--$76\times$ faster),
  near-zero RMSE is prioritised, or $\sigma>0$ conditions are expected.
\end{itemize}

%% ============================================================
\section{Root-Cause Analysis: Newton's Gravity Measurement Bug}
\label{sec:bugfix}

\subsection{Symptom}

M4 reported $\Rsq \in [0.42, 0.87]$ on Newton's gravitational force across
all 11 tested conditions in v1, despite the LLM generating the correct formula
on every run.

\subsection{Root Cause}

\texttt{evaluate\_llm\_formula} used:
\begin{verbatim}
r2 = 1 - (ss_res / ss_tot) if ss_tot > 1e-10 else 0
\end{verbatim}
For Newton's gravity at tested parameter ranges, outputs are $\sim 10^{-11}$\,N,
giving $\mathrm{ss\_tot} \approx 1.8 \times 10^{-18}$ — eight orders of
magnitude below the threshold.  The function always returned
\texttt{r2 = 0.0}, triggering NN fallback.

\subsection{Fix Applied}

\begin{verbatim}
_tol = 1e-10 * float(np.max(np.abs(y_true))**2) * len(y_true)
r2   = 1 - (ss_res / ss_tot)
        if ss_tot > max(_tol, 1e-300) else 0
\end{verbatim}

The same fix was applied at three code locations:
\texttt{evaluate\_llm\_formula} (line 462),
\texttt{train\_nn} (line 384), and
the ensemble blend $\Rsq$ calculation (line 700).

\subsection{Additional Fixes in the Same Commit}

\begin{itemize}
  \item \texttt{force\_llm} flag: now correctly bypasses NN training.
  \item \texttt{PureLLMBaseline}: validates non-empty \texttt{python\_code}.
  \item Training epochs: corrected from 300 to 1000 (scheduler requires $\geq 1000$).
  \item Log-space guard: unconditionally enables log-space training when
    \texttt{y\_range > 3.0} decades.
\end{itemize}

\paragraph{General contribution.}
Any SR benchmark harness using absolute sum-of-squares thresholds will
silently misclassify correct formulas on equations with output values far
from unity.  The relative threshold fix should be adopted as standard practice.

%% ============================================================
\section{Discussion}
\label{suppb:discussion}

\subsection{Hardcoding and LLM Benchmark Integrity}

Of \PureLLM{}'s 30 noiseless passes, 11 (37\%) are attributable to
memorisation, inflating the reported pass rate from 63\% (genuine) to
100\% (inclusive).  Explicit hardcoding audits should be standard in any
benchmark involving LLM-based SR.

\subsection{M3 vs.\ M4 Trade-off}

With the measurement bug resolved, M3 and M4 are \textbf{equivalent on
recovery rate and $\Rsq$}: both achieve 100\% recovery and median
$\Rsq = 1.000000$ at every tested noisy condition.  The distinction is:
M3 is interpretability-optimised (exact symbolic formulas); M4 is
efficiency-optimised ($1.5$--$76\times$ faster, near-zero RMSE).

The 66\% tie rate across 330 equation-condition pairs quantifies their
practical equivalence.

%% ============================================================
\section{Limitations}
\label{suppb:limitations}

\begin{enumerate}
  \item \textbf{Benchmark scope.} 30 equations, 2--5 input variables.
    Higher-dimensional and PDE targets remain untested.
  \item \textbf{Memorisation scope.} The hardcoding audit detects verbatim
    recall but not partial memorisation.
  \item \textbf{Baseline breadth.} Only M3 and M4 covered by the sweep;
    adding PySR and DSR sweeps would strengthen generalisation claims.
  \item \textbf{M3 noiseless timing.} The 841\,s/eq cost is an offline
    exhaustive-search cost, not inference time.
  \item \textbf{No formula complexity metric.} Expression tree depth and
    operator count should accompany $\Rsq$ in future evaluations.
  \item \textbf{[Open, process.]} The project's CI dashboard (July~9 run)
    should not be cited: its own upstream notebooks (NB-01--NB-06) report
    ``not found'' for that run, while its summary asserts ``0 open
    issues'' and specific claims resolved. \texttt{build\_report.py}
    needs to hard-fail on missing notebook output rather than emit a
    stale green summary, and the dashboard needs a genuine re-run before
    any of its figures are used.
  \item \textbf{[Resolved.]} The legacy exp2 baseline-lock threshold
    fix ($R^2 \geq 0.999999$ replacing $R^2 \geq 0.9999$, which had
    inflated 74/90 to 81/90) is now live-confirmed at 74/90 against a
    fresh CI run, and reproduced independently a second time via a
    patched Gate~C schema check; 81/90 should not be cited as current.
\end{enumerate}

%% ============================================================
\section{JMLR Readiness Assessment}
\label{sec:jmlr}

\noindent\textbf{Note (v3.0 update, April 2026).}
The readiness metrics in Table~\ref{tab:jmlr} below are based on the v2
benchmark run (30 Feynman equations, noiseless/noise sweep protocol).
The main paper (\texttt{jmlr\_paper\_main.tex}) now reports v3.0
results on 74 DeFi tasks under the PCA-directed aggressive extrapolation
split: HypatiaX achieves 89.2\,\% near-perfect success, eliminating all 6
LLM catastrophic failures.  These are a \emph{different and more stringent}
evaluation than the noiseless recovery metrics here; the v2 Feynman readiness
assessment remains valid for this document's scope.

\begin{table}[ht]
\centering
\caption{JMLR readiness assessment (v2 Feynman results, post bug-fix).
  See note above for v3.0 DeFi results reported in the main paper.}
\label{tab:jmlr}
\small
\renewcommand{\arraystretch}{1.3}
\begin{tabular}{l C{1.5cm} C{1.5cm} L{5.0cm}}
\toprule
\textbf{Criterion} & \textbf{M3} & \textbf{M4} & \textbf{Required action}\\
\midrule
Accuracy ($\Rsq$)       & \textcolor{passgreen}{\textbf{PASS}} & \textcolor{passgreen}{\textbf{PASS}} &
  Both: median $\Rsq=1.000000$\\
Noise robustness        & \textcolor{passgreen}{\textbf{STRONG}} & \textcolor{passgreen}{\textbf{STRONG}} &
  M3 100\% all $\sigma$; M4 100\% at $\sigma>0$\\
Data efficiency         & \textcolor{passgreen}{\textbf{STRONG}} & \textcolor{passgreen}{\textbf{STRONG}} &
  Convergence at $n=50$ for both\\
Catastrophic failures   & \textcolor{passgreen}{\textbf{CLEAN}} & \textcolor{passgreen}{\textbf{CLEAN}} &
  0/30 for both; document Newton bug\\
Computational cost      & \textcolor{warnAmber}{\textbf{COND.}} & \textcolor{passgreen}{\textbf{STRONG}} &
  M3: frame 841\,s/eq as offline search\\
Statistical significance & \textcolor{passgreen}{\textbf{PASS}} & \textcolor{passgreen}{\textbf{PASS}} &
  Add Wilcoxon test; report 66\% tie rate\\
Hardcoding disclosure   & \multicolumn{2}{c}{\textcolor{warnAmber}{\textbf{REQUIRED}}} &
  \PureLLM: disclose 11/30 hardcoded\\
Process/CI credibility  & \multicolumn{2}{c}{\textcolor{warnAmber}{\textbf{OPEN}}} &
  Baseline-lock now live-confirmed (74/90); do not cite July~9 dashboard until re-run (\S\ref{suppb:limitations})\\
\midrule
\textbf{Overall verdict} & \textcolor{passgreen}{\textbf{READY}} & \textcolor{passgreen}{\textbf{READY}} &
  Ready on accuracy/robustness grounds; one process item above remains open\\
\bottomrule
\end{tabular}
\end{table}

%% ============================================================
\section{Conclusions}
\label{sec:conclusions}

\begin{enumerate}
  \item \textbf{Both M3 and M4 achieve 100\% recovery} at all tested noisy
    conditions ($0.5$--$10\%$ noise) and all sample sizes ($50$--$1000$).
  \item \textbf{M3 outperforms M4 at $\sigma=0$}: 100\% vs 90\% noiseless;
    three further M4 noiseless failures remain (Lorentz force, Photon energy,
    Zeeman energy).
  \item \textbf{M4 is $1.5$--$76\times$ faster} with near-zero RMSE; for
    production sweeps its throughput advantage is decisive.
  \item \textbf{Both substantially exceed published results on the 30-equation subset}:
    \EHD{} at 96.7\% vs.\ AI~Feynman~2.0 at 79.3\% ($+17.4$\,pp); \HDS{}
    at 90.0\% ($+10.7$\,pp).  Published baselines were evaluated on all 100
    Feynman equations under SRBench; these are subset-specific comparisons.
  \item \textbf{The Newton failure in v1 was a measurement bug}: absolute
    $\Rsq$ threshold misclassified correct formulas for tiny-scale equations.
    The relative threshold fix is a general contribution to SR benchmarking.
  \item \textbf{Four methodological issues corrected}: undisclosed hardcoding,
    scale-biased \rmse, missing significance tests, mismatched thresholds.
\end{enumerate}

%% ── Bibliography ─────────────────────────────────────────────────────────────
\bibliography{references}

%% ============================================================
\appendix

\section{Reproducibility Figure Generation Inventory}
\label{app:figures}

\noindent This appendix documents the auxiliary figures produced by
\texttt{generate\_figures.py} for reproducibility purposes. It is
distinct from the figures actually embedded in the compiled
manuscript (see \texttt{figures\_inventory.tex}, 5 files, listed in
\texttt{jmlr\_paper\_main.tex}).

\begin{table}[H]
\centering
\caption{Complete figure inventory; all files in \texttt{figures/}.
  All 13 generated figures must be regenerated from v2 result JSONs.}
\label{tab:figlist}
\small
\renewcommand{\arraystretch}{1.25}
\begin{tabular}{L{4.8cm} L{1.5cm} L{1.8cm} L{3.2cm}}
\toprule
\textbf{Filename} & \textbf{Sweep} & \textbf{Section} & \textbf{Content}\\
\midrule
\texttt{fig1\_r2\_vs\_noise.pdf}       & Noise & \ref{sec:noise}        & Median $\Rsq$ vs $\sigma$\\
\texttt{fig2\_rmse\_vs\_noise.pdf}     & Noise & \ref{sec:noise}        & Median RMSE vs $\sigma$\\
\texttt{fig3\_time\_vs\_noise.pdf}     & Noise & \ref{sec:noise}        & Avg time vs $\sigma$\\
\texttt{fig4\_r2\_vs\_n.pdf}           & SC    & \ref{sec:sc}           & Median $\Rsq$ vs $n$\\
\texttt{fig5\_rmse\_vs\_n.pdf}         & SC    & \ref{sec:sc}           & Median RMSE vs $n$\\
\texttt{fig6\_time\_vs\_n.pdf}         & SC    & \ref{sec:sc}           & Median time vs $n$\\
\texttt{fig7\_recovery\_vs\_noise.pdf} & Noise & \ref{sec:noise}        & Recovery rate vs $\sigma$\\
\texttt{fig8\_recovery\_vs\_n.pdf}     & SC    & \ref{sec:sc}           & Recovery rate vs $n$\\
\texttt{fig9\_minr2\_vs\_noise.pdf}    & Noise & \ref{sec:noise}        & Min $\Rsq$ vs $\sigma$\\
\texttt{fig10\_r2\_boxplot\_noise.pdf} & Noise & \ref{sec:convergence}  & Per-eq $\Rsq$ box plots\\
\texttt{fig11\_recovery\_heatmap.pdf}  & Both  & \ref{sec:convergence}  & Recovery heatmap $\sigma \times n$\\
\midrule
\texttt{fig\_runtime\_comparison.png}  & ---   & \ref{sec:noiseless}    & Runtime bar chart, 6 methods\\
\texttt{fig\_comparative\_table.png}   & ---   & \ref{sec:noiseless}    & Domain $\times$ method table\\
\bottomrule
\end{tabular}
\end{table}

\noindent\textbf{Generation note:} Regenerate from
\texttt{noise\_sweep\_20260316\_192711.json} and
\texttt{sample\_complexity\_20260316\_193447.json} (v2 corrected
harness; see Section~\ref{sec:r2_measurement_bugfix} for v1 vs.\ v2 provenance).

\noindent\textbf{Relationship to the submitted PDF:} the 13 files
above are reproducibility documentation only and are not
\verb|\includegraphics|'d anywhere in this document or in the main
manuscript. The compiled manuscript embeds a separate, disjoint set
of 5 figures directly via \verb|\includegraphics| (see
\texttt{jmlr\_paper\_main.tex}). Only those 5 files are bundled in
\texttt{submission/sources/figures/}; the 13 listed here should be
regenerated from source by anyone who wants them, not shipped in the
submission package.

\section{Method Abbreviations}
\label{app:abbrev}

\begin{tabular}{L{1.8cm} L{2.5cm} L{6.0cm}}
\toprule
\textbf{Short} & \textbf{Code} & \textbf{Full name}\\
\midrule
\PureLLM & M1 & PureLLM Baseline (core)\\
\INN     & M2 & ImprovedNN (core)\\
\EHD     & M3 & EnhancedHybridSystemDeFi (core)\\
\HSL     & M4 & HybridSystemLLMNN all-domains (core)\\
\SEL     & M5 & SymbolicEngineWithLLM (tools)\\
\HDS     & M6 & HybridDiscoverySystem v50\_2 (tools)\\
\bottomrule
\end{tabular}

\section{PureLLM Required Disclosure}
\label{app:disclosure}

Required in any JMLR submission including \PureLLM{} results:

\begin{quote}
\textit{The PureLLM Baseline returns pre-stored formulas for 11 of 30
benchmark equations (\texttt{is\_hardcoded: true}); its reported pass rate
of 100\% includes these cases and does not solely reflect numerical
regression from data.  The genuine-inference pass rate (19 non-hardcoded
equations) is 100\%.}
\end{quote}

% ====================================================================
% NEW APPENDICES (extracted from jmlr_paper_main.tex)
% ====================================================================
\section{Feynman Benchmark Reproducibility Statement}
\label{app:repro_statement}

\subsection*{Code, Data, and Protocol}

All experimental results in the main paper's \emph{Benchmark Design} section were produced
using \texttt{experiment\_protocol\_benchmark\_v2.py} (version~2.0,
2026-03-02), which will be released alongside the final paper.
The protocol file encodes all data-generation functions, random seeds, noise
levels, and recovery thresholds as class constants and is self-documenting:

\begin{verbatim}
python experiment_protocol_benchmark_v2.py --describe
\end{verbatim}

\noindent The 30-equation Feynman subset is defined in
\texttt{\_build\_feynman\_equations()} (lines 480--953 of the protocol file).
Each equation specifies its Feynman ID, ground-truth formula, variable
ranges, and generating function; data are produced analytically, eliminating
external dataset dependencies for the physics equations.

\subsection*{Exact Reproduction Commands}

\paragraph{Phase~2 (noisy, $R^2>0.995$).\textbf{[EDITORIAL NOTE: no table in this
document reports a noisy-protocol, six-method, 30-equation aggregate at the
0.995 threshold --- the original \texttt{tab:feynman30} label does not exist
anywhere in the paper or supplements. Please confirm whether this should
point to a specific table (e.g.\ a per-equation noisy breakdown analogous to
Table~\ref{tab:overall}) or be removed.]}}
\begin{verbatim}
python run_comparative_suite_benchmark_v2.py \
    --methods 1 --samples 200 --no-llm-cache
\end{verbatim}
Output: \texttt{protocol\_core\_noisy\_<timestamp>.json}.
The \texttt{--no-llm-cache} flag forces independent fresh API calls for all
30 Pure~LLM evaluations; this was verified via cache-clearing.

\paragraph{Phase~3 (noiseless, $R^2>0.9999$, Table~\ref{tab:overall}).}
\begin{verbatim}
python run_comparative_suite_benchmark_v2.py \
    --noiseless --threshold 0.9999       \
    --nn-seeds 3 --samples 200           \
    --method-timeout 900                 \
    --pysr-timeout 900
\end{verbatim}
Output: \texttt{protocol\_core\_noiseless\_20260304\_154510.json}
(full run: 30/30 equations, all 6 methods, $\approx$1\,h\,13\,min wall-clock
on Intel Celeron T3100, 2~cores; faster on modern multi-core hardware).

\textbf{Note on timeout upgrade:} \texttt{--method-timeout} was set to 300\,s
for tests 1--18 and upgraded to 900\,s for tests 19--30. Arrhenius (test~4)
ran under the 300\,s limit and hung for 27574\,s before the Julia process was
killed; this is recorded as a genuine failure independent of the v2 bug fix.
Nernst (test~6) also timed out under the 300\,s default; counted as failure.

\subsection*{Random Seeds}

\begin{itemize}
  \item Base seed: \texttt{42} (class constant
    \texttt{BenchmarkProtocol.seed = 42}, line~1091 of the protocol file).
    Matches the Core~15 benchmark seed for cross-campaign consistency.
  \item Per-equation seed: \texttt{seed + offset $\times$ 7}, where
    \texttt{offset} is the 0-based index in the equation list
    (line~1208: \texttt{eq.generate(n, noise, seed=s + offset * 7)}).
    This ensures each of the 30 equations receives a distinct, reproducible
    seed regardless of call order.
  \item Neural network seeds (Phase~3): 3 independent seeds per equation;
    median $R^2$ reported. Exact seed values are logged in the output JSON.
\end{itemize}

\subsection*{Software Environment}

\begin{table}[H]
\centering
\caption{Software and hardware environment for Feynman benchmark runs
         (verified 18 March 2026 via \texttt{hardware\_info.py} and
         \texttt{pip show}).}
\label{tab:software_env}
\small
\begin{tabular}{ll}
\toprule
\textbf{Component} & \textbf{Version / details} \\
\midrule
Python           & 3.12.2 \\
PySR             & 1.5.9 \\
Julia runtime    & 1.12.2 \\
juliacall        & 0.9.26 (Python$\leftrightarrow$Julia bridge used by PySR) \\
NumPy            & 2.3.5 \\
SciPy            & 1.16.3 \\
SymPy            & 1.14.0 \\
Pandas           & 2.3.3 \\
Matplotlib       & 3.10.7 \\
Seaborn          & 0.13.2 \\
scikit-learn     & 1.7.2 \\
Anthropic SDK    & 0.73.0 (\texttt{claude-sonnet-4-20250514},
                   temperature~0.0) \\
\midrule
Hardware         & Intel Celeron T3100 @ 1.90\,GHz (2~cores);
                   3\,GB DDR2; no GPU \\
OS               & Kali GNU/Linux Rolling \\
PySR populations & 2 (hardware limit: 2-core CPU) \\
\bottomrule
\end{tabular}
\end{table}

\noindent\textbf{Note on Anthropic SDK version compatibility.}
\texttt{anthropic==0.73.0} is substantially newer than early-2024 releases.
The Messages API interface is stable across 0.18$+$ versions; earlier
installations may require minor call-site adjustments but the prompts and
response parsing are otherwise unchanged.

\noindent PySR uses Julia under the hood; the 900\,s
\texttt{--pysr-timeout} was required to accommodate Julia JIT
initialisation on the first run. Subsequent equations in the same process
are faster; the 27574\,s Arrhenius hang is specific to single-variable
flat-landscape equations where PySR's complexity search does not converge
within the timeout.

\subsection*{Measurement Correction (v1\,$\to$\,v2)}
\label{sec:r2_measurement_bugfix}

A bug was identified on 2026-03-16 in \texttt{evaluate\_llm\_formula}
and corrected before final submission. The function used a hardcoded
absolute sum-of-squares guard (\texttt{ss\_tot > 1e-10}) that silently
returned $R^2=0$ for equations whose output magnitudes are far below unity
(Newton gravity $\sim\!10^{-11}$\,N, Zeeman splittings $\sim\!10^{-23}$\,J,
Coulomb forces at $\mu$C scale). The fix scales the threshold relative
to the data:

\begin{verbatim}
# v1 (buggy):
r2 = 1 - (ss_res / ss_tot) if ss_tot > 1e-10 else 0

# v2 (fixed):
_tol = 1e-10 * float(np.max(np.abs(y_true))**2) * len(y_true)
r2 = 1 - (ss_res / ss_tot) if ss_tot > max(_tol, 1e-300) else 0
\end{verbatim}

The same fix was applied at three locations:
\begin{enumerate}
  \item \texttt{evaluate\_llm\_formula} (line 462) — primary failure site
  \item \texttt{train\_nn} (line 384) — NN self-evaluation
  \item Ensemble-blend $R^2$ calculation (line 700) — candidate scoring
\end{enumerate}

Additional fixes in the same commit:
\begin{itemize}
  \item \texttt{force\_llm} flag: \texttt{hybrid\_predict} now correctly
    bypasses NN training when the LLM formula passes the threshold.
  \item \texttt{PureLLMBaseline} passthrough: \texttt{generate\_llm\_formula}
    now validates non-empty \texttt{python\_code} before accepting.
  \item Training epochs: corrected from 300 to 1000 to satisfy the learning
    rate scheduler's warm-restart requirement ($\geq\!1000$ epochs).
  \item Log-space guard: \texttt{\_is\_pole} heuristic now unconditionally
    enables log-space training when \texttt{y\_range > 3.0} decades,
    fixing suppression on wide-range power-law equations.
\end{itemize}

The \emph{Summary of Changes} table in Supplementary~A (\emph{Routing Improvements for Hybrid LLM--Neural-Network Systems}), and Section~\ref{sec:r2_measurement_bugfix} of this document, summarise all
aggregate recovery changes. All results reported in the paper use the
v2 corrected harness. The v1 JSON output files are retained for audit
purposes and are available upon request.

\subsection*{Per-Equation v2 Data}

The complete per-equation $R^2$ values under the v2 corrected harness are
stored in:
\begin{itemize}
  \item \texttt{protocol\_core\_noiseless\_20260304\_154510.json} —
    Phase~3 noiseless, all 6 methods, 30 equations.
  \item \texttt{noise\_sweep\_20260316\_192711.json} —
    M3/M4 noise sweep ($\sigma \in \{0, 0.5, 1, 5, 10\}\%$, v2 corrected).
  \item \texttt{sample\_complexity\_20260316\_193447.json} —
    M3/M4 sample complexity sweep ($n \in \{50, 100, 200, 500, 750, 1000\}$,
    v2 corrected).
\end{itemize}
These files will be released as supplementary material. The v1 superseded
files (\texttt{noise\_sweep\_20260315\_091018.json},
\texttt{sample\_complexity\_20260315\_124310.json}) should not be used
for the final submission figures.

\subsection*{PureLLM Disclosure}

Following~\citet{memorization2024audit}, each PureLLM noiseless pass
was audited for memorisation. \textbf{11 of 30 PureLLM passes are
flagged \texttt{is\_hardcoded=True}}: Michaelis--Menten, Logistic growth,
Allometric scaling, Nernst, Clausius--Mossotti, Lorentz force, Gaussian PDF,
Zeeman energy, Bose--Einstein, Planck radiance, and Fourier heat. These
equations return the ground-truth formula verbatim with zero cross-seed
variation, consistent with training-data recall rather than inference.
The genuine-inference pass rate (19 non-hardcoded equations) is 100\%.
PureLLM is therefore listed for completeness in all tables but excluded
from all scientific comparisons (the main paper's Results section).

A controlled re-run with \texttt{--no-llm-cache} confirmed that all
30 equations issued independent API calls and produced identical scores
(median $R^2=0.9976$, std $=2.5\times10^{-4}$), ruling out any
caching artefact.

\subsection*{Citation Key Corrections}

The following citation key conflicts have been resolved in \texttt{references.bib}
(updated 2026-06):

\begin{enumerate}
  \item \texttt{cranmer2023interpretable} renamed to \texttt{cranmer2023pysr}
        (canonical key matching all \texttt{\textbackslash cite} calls in the main paper;
        same paper arXiv:2305.01582). \textbf{Resolved.}
  \item \texttt{udrescu2020aifeynman} duplicate removed; \texttt{udrescu2020ai}
        retained as the single canonical key in \texttt{references.bib}.
        This supplementary document uses \texttt{udrescu2020aifeynman} in its own
        inline \texttt{thebibliography} and remains self-consistent. \textbf{Resolved.}
  \item \texttt{meidani2023snip} (arXiv preprint key) replaced by
        \texttt{meidani2024snip} (published ICLR~2024 key) in
        \texttt{references.bib}. \textbf{Resolved.}
  \item \texttt{lacava2021contemporary} entry confirmed present in
        \texttt{references.bib}. \textbf{No action needed.}
\end{enumerate}

\subsection*{SNIP Comparison Note}

SNIP~\citep{meidani2024snip} does \emph{not} report recovery rates.
It evaluates symbolic regression using a Pareto front of median
$R^2$ vs.\ equation complexity on 119 Feynman equations, reporting
median $R^2 = 0.882$ (vs.\ AI~Feynman baseline: $0.798$).
Placing SNIP's median $R^2$ alongside HypatiaX's recovery rate
(\% at $R^2>0.9999$) would be a category error.
The only valid comparison uses the shared median $R^2$ axis:
HypatiaX Hybrid DeFi achieves median $R^2 = 1.0000$ on our 30-equation
subset vs.\ SNIP's $0.882$ on the full 119-equation SRBench set
(different subsets; comparison is indicative only).
The Nguyen-12 benchmark table in the main paper provides this comparison.

% ============================================================================
% APPENDIX A: COMPLETE PROOF OF THEOREM
% ============================================================================


\section{Hyperparameter Settings}
\label{app:hyperparameters}

\begin{table}[H]
\centering
\caption{PySR Configuration for Symbolic Regression}
\label{tab:hyperparameters}
\begin{tabular}{@{}lc@{}}
\toprule
\textbf{Parameter} & \textbf{Value} \\
\midrule
Population size & 100 \\
Generations & 500 \\
Populations (parallel) & 2 (hardware limit: 2-core CPU) \\
Crossover probability & 0.8 \\
Mutation probability & 0.1 \\
Complexity penalty & 0.01 \\
Tournament size & 3 \\
Elitism & 5 \\
Max expression depth & 15 \\
Optimization iterations & 50 \\
\midrule
\textbf{Validation Thresholds} & \\
\midrule
Acceptance $R^2$ (Criterion 1) & 0.99 \\
Validation score (Criterion 1) & 30 \\
Acceptance $R^2$ (Criterion 2) & 0.95 \\
Validation score (Criterion 2) & 80 \\
Extrapolation $R^2$ & 0.80 \\
Numerical stability ($\alpha$) & 0.95 \\
\bottomrule
\end{tabular}
\end{table}

\begin{table}[H]
\centering
\caption{Complete Hyperparameter Specifications for All Methods}
\label{tab:hyperparameters_complete}
\small
\begin{tabular}{lll}
\toprule
\textbf{Method} & \textbf{Parameter} & \textbf{Value} \\
\midrule
\multirow{8}{*}{Neural Network} 
& Architecture & [64, 32, 1] \\
& Activation & ReLU \\
& Optimizer & Adam \\
& Learning rate & 0.01 \\
& Batch size & Full batch (160 samples) \\
& Epochs & 200 \\
& Weight initialization & Xavier uniform \\
& Dropout & None \\
\midrule
\multirow{7}{*}{Pure LLM}
& Model & claude-sonnet-4-20250514 \\
& Temperature & 0.0 (deterministic) \\
& Max tokens & 1024 \\
& Top-p & 1.0 \\
& Few-shot examples & 2-3 per domain \\
& Prompt template & Domain-specific \\
& Retry on failure & 3 attempts \\
\midrule
\multirow{12}{*}{Hybrid v50\_2}
& SR iterations & 50 \\
& SR populations & 2 (hardware limit: 2-core CPU) \\
& Binary operators & [+, -, *, /] \\
& Unary operators & [exp, log, sqrt, sin, cos] \\
& Parsimony coefficient & 0.01 \\
& Tournament size & 10 \\
& Mutation rate & 0.1 \\
& Crossover rate & 0.9 \\
& LLM candidates & 5 \\
& LLM temperature & 0.3 \\
& Max complexity & 30 \\
& Validation threshold ($R^2$) & 0.99 \\
\bottomrule
\end{tabular}
\end{table}

% ============================================================================
% APPENDIX C: DETAILED EXPERIMENTAL RESULTS
% ============================================================================


\section{Detailed Experimental Results}
\label{app:detailed_results}

\subsection{Pure LLM Baseline (40 tests)}

\begin{table}[H]
\centering
\caption{Pure LLM Performance by Test Case}
\label{tab:llm_detailed}
\begin{small}
\begin{tabular}{@{}llccp{5.2cm}@{}}
\toprule
\textbf{Domain} & \textbf{Test} & \textbf{$R^2$} & \textbf{Result} & \textbf{Failure Mode} \\
\midrule
\multicolumn{5}{c}{\textit{Classical Scientific Domains (20 tests)}} \\
\midrule
Materials & Stress-strain & 1.00 & Pass & --- \\
Materials & Young's modulus & 1.00 & Pass & --- \\
Materials & Poisson ratio & 1.00 & Pass & --- \\
Materials & Thermal expansion & 1.00 & Pass & --- \\
Fluid & Reynolds number & 1.00 & Pass & --- \\
Fluid & Bernoulli equation & 1.00 & Pass & --- \\
Fluid & Drag coefficient & 1.00 & Pass & --- \\
Fluid & Flow rate & 1.00 & Pass & --- \\
Thermo & Ideal gas law & 1.00 & Pass & --- \\
Thermo & Heat capacity & 1.00 & Pass & --- \\
Thermo & Carnot efficiency & 1.00 & Pass & --- \\
Thermo & Stefan-Boltzmann & 1.00 & Pass & --- \\
Mechanics & Newton's 2nd law & 1.00 & Pass & --- \\
Mechanics & Kinetic energy & 1.00 & Pass & --- \\
Mechanics & Gravitational force & 1.00 & Pass & --- \\
Mechanics & Hooke's law & 1.00 & Pass & --- \\
Chemistry & Arrhenius equation & 1.00 & Pass & --- \\
Chemistry & pH calculation & 1.00 & Pass & --- \\
Chemistry & Nernst equation & 0.50 & Fail & Wrong sign convention \\
Chemistry & Rate law & 1.00 & Pass & --- \\
\midrule
\multicolumn{5}{c}{\textit{Decentralized Finance Domains (20 tests)}} \\
\midrule
AMM & Constant product & 1.00 & Pass & --- \\
AMM & Price impact & 0.98 & Pass & --- \\
AMM & Impermanent loss & $-1.26$ & Fail & Unit inconsistency (missing $\times$100) \\
AMM & Swap fee & 0.82 & Pass & --- \\
VaR & Parametric VaR 95\% & $-10.05$ & Fail & Wrong sign (z-score) \\
VaR & Parametric VaR 99\% & 1.00 & Pass & --- \\
VaR & Historical VaR & 0.95 & Pass & --- \\
VaR & Modified VaR & $-2.15$ & Fail & Scipy API misuse \\
Liquidity & APY calculation & 1.00 & Pass & --- \\
Liquidity & Capital efficiency & 0.87 & Pass & --- \\
Liquidity & Fee earnings & --- & Fail & Non-executable (tutorial mode) \\
Liquidity & Utilization ratio & --- & Fail & Non-executable (tutorial mode) \\
ES & Expected Shortfall 95\% & $-4.12$ & Fail & Scipy \texttt{norm.pdf} returns object \\
ES & Expected Shortfall 99\% & $-3.87$ & Fail & Distribution tail error \\
ES & Conditional VaR & 0.88 & Pass & --- \\
ES & Tail risk & $-1.92$ & Fail & Statistical definition error \\
Liquidation & Long position & --- & Fail & Tutorial text only \\
Liquidation & Short position & --- & Fail & Tutorial text only \\
Liquidation & Margin call & --- & Fail & Incomplete derivation \\
Liquidation & Leverage ratio & --- & Fail & Multi-step missing \\
\midrule
\multicolumn{3}{l}{\textbf{Classical Average:}} & \textbf{19/20 (95\%)} & \\
\multicolumn{3}{l}{\textbf{DeFi Average:}} & \textbf{8/20 (40\%)} & \\
\multicolumn{3}{l}{\textbf{Overall:}} & \textbf{27/40 (67.5\%)} & \\
\bottomrule
\end{tabular}
\end{small}
\end{table}

\paragraph{Key Findings:} Pure LLM performance strongly correlates with equation prevalence in training data (Spearman $\rho = 0.89$, $p < 0.001$)\citep{kambhampati2024llms}. Classical equations achieve 95\% success rate while specialised domains (DeFi) drop to 40\%, consistent with the reproductive behaviour characterisation in the main paper's \emph{LLM Reproductive Behaviour} finding.

\subsection{DeFi-Optimized Results (23 tests)}

\begin{table}[H]
\centering
\caption{DeFi-Optimized HypatiaX Performance}
\label{tab:defi_detailed}
\begin{small}
\resizebox{\textwidth}{!}{%
\begin{tabular}{@{}lccccc@{}}
\toprule
\textbf{Test} & \textbf{$R^2$} & \textbf{Validation} & \textbf{Time (s)} & \textbf{Extrap.} & \textbf{Result} \\
\midrule
Constant product (xy=k) & 1.0000 & 100.0 & 45 & Pass & Pass \\
Price impact & 0.9998 & 95.2 & 120 & Pass & Pass \\
Impermanent loss & 0.9995 & 92.1 & 180 & Pass & Pass \\
Swap fee calculation & 1.0000 & 98.5 & 60 & Pass & Pass \\
VaR 95\% parametric & 0.9988 & 88.3 & 150 & Pass & Pass \\
VaR 99\% parametric & 0.9990 & 90.1 & 145 & Pass & Pass \\
Historical VaR & 0.9985 & 87.5 & 135 & Pass & Pass \\
Modified VaR & 0.9982 & 85.9 & 160 & Pass & Pass \\
APY simple & 1.0000 & 100.0 & 50 & Pass & Pass \\
APY compound & 0.9995 & 94.8 & 110 & Pass & Pass \\
Capital efficiency & 0.9940 & 78.5 & 240 & Pass & Pass \\
Fee earnings & 0.9998 & 96.2 & 90 & Pass & Pass \\
Expected Shortfall 95\% & 0.9987 & 89.7 & 170 & Pass & Pass \\
Expected Shortfall 99\% & 0.9985 & 88.2 & 175 & Pass & Pass \\
Conditional VaR & 0.9990 & 91.4 & 155 & Pass & Pass \\
Tail risk metric & 0.9983 & 86.8 & 165 & Pass & Pass \\
Liquidation (long) & 0.9992 & 93.6 & 125 & --- & Pass \\
Liquidation (short) & 0.9994 & 94.1 & 120 & --- & Pass \\
Margin call threshold & 0.9996 & 95.8 & 105 & --- & Pass \\
Effective leverage & 0.9998 & 97.2 & 85 & --- & Pass \\
Staking rewards & 1.0000 & 100.0 & 55 & --- & Pass \\
Dilution effect & 0.9993 & 93.8 & 115 & --- & Pass \\
Kelly criterion & 0.9988 & 89.5 & 195 & Pass & Pass \\
\midrule
\textbf{Average} & \textbf{0.9990} & \textbf{91.3} & \textbf{115.2} & \textbf{16/16} & \textbf{23/23} \\
\bottomrule
\end{tabular}%
}% end resizebox
\end{small}
\end{table}

\paragraph{Extrapolation Tests:} Sixteen tests designated for extrapolation validation all passed with average extrapolation $R^2 = 0.9650$, indicating reliable generalization to the tested novel regimes on this benchmark\citep{balestriero2021high,neyshabur2017exploring,zhang2021understanding}.

\subsection{Complete Test Suite Details}

\begin{table}[H]
\centering
\caption{Ground Truth Formulas with Exact Constants}
\label{tab:formulas_detailed}
\scriptsize
\begin{tabular}{lll}
\toprule
\textbf{Equation} & \textbf{Formula} & \textbf{Constants} \\
\midrule
Arrhenius & $k = A\exp(-E_a/(RT))$ & $A=10^{11}$, $E_a=80000$ J/mol, $R=8.314$ J/(mol$\cdot$K) \\
Henderson-Hasselbalch & $\text{pH} = \text{p}K_a + \log_{10}([A^-]/[\text{HA}])$ & $\text{p}K_a = 6.5$ \\
Rate Law & $r = k[A]^2[B]$ & $k = 0.5$ M$^{-2}$s$^{-1}$ \\
Allometric Scaling & $Y = aM^b$ & $a = 3.5$, $b = 0.75$ \\
Michaelis-Menten & $v = V_{\max}[S]/(K_m + [S])$ & $V_{\max} = 50$ $\mu$M/s, $K_m = 10$ $\mu$M \\
Logistic Growth & $dN/dt = rN(1-N/K)$ & $r = 0.3$ day$^{-1}$, $K = 1000$ \\
Kinetic Energy & $E = \frac{1}{2}mv^2$ & Dimensionless \\
Gravitational Force & $F = Gm_1m_2/r^2$ & $G = 6.674 \times 10^{-11}$ N$\cdot$m$^2$/kg$^2$ \\
Ideal Gas Law & $P = nRT/V$ & $R = 8.314$ J/(mol$\cdot$K) \\
Impermanent Loss & $\text{IL} = 2\sqrt{r}/(1+r) - 1$ & Dimensionless \\
Price Impact & $\Delta p = \Delta x/(x + \Delta x)$ & Dimensionless \\
Constant Product & $y = k/x$ & $k = 10^6$ \\
VaR 95\% & $\text{VaR} = P\sigma \cdot z_{0.95}$ & $z_{0.95} = 1.645$ \\
Liquidation Price & $p_{\text{liq}} = p_0(1 - 1/(L \cdot m))$ & $m = 0.8$ (maintenance margin) \\
Portfolio Variance & $\sigma_p = \sqrt{\sigma_1^2 + \sigma_2^2 + 2\rho\sigma_1\sigma_2}$ & Dimensionless \\
\bottomrule
\end{tabular}
\end{table}

% ============================================================================
% APPENDIX D: STATISTICAL TEST DETAILS (CORRECTED VALUES)
% ============================================================================


\section{Statistical Test Details}
\label{app:statistical_tests}

\subsection{Mann-Whitney U Test for Extrapolation Performance}

For medium extrapolation regime (2$\times$ training range):

\paragraph{Test Setup:}
\begin{itemize}
\item \textbf{Sample sizes}: $n_{\text{Hybrid}} = 14$, $n_{\text{NN}} = 13$ (CORRECTED)
\item \textbf{Null hypothesis}: $H_0: E_{\text{hybrid}} \geq E_{\text{NN}}$ (Hybrid is not better)
\item \textbf{Alternative hypothesis}: $H_1: E_{\text{hybrid}} < E_{\text{NN}}$ (Hybrid is better)
\item \textbf{Significance level}: $\alpha = 0.05$ (one-tailed test)
\end{itemize}

\paragraph{Descriptive Statistics:}
\begin{itemize}
\item \textbf{Hybrid v50\_2}: Mean = 0.0\%, Median = 0.0\%, SD = 0.0\%, Range = [0.0, 0.0]
\item \textbf{Neural Network}: Mean = 1231\%, Median = 86.7\%, SD = 1842\%, Range = [12.3, 5847\%] (CORRECTED)
\end{itemize}

\paragraph{Test Results:}
\begin{itemize}
\item \textbf{Test statistic}: $U = 0$ (all Hybrid errors < all NN errors)
\item \textbf{P-value}: $p = 1.11 \times 10^{-6}$ (CORRECTED)
\item \textbf{Critical value}: $U_{0.05} = 45$ for one-tailed test with n=14 vs n=13
\item \textbf{Conclusion}: $U < U_{0.05}$ $\Rightarrow$ Reject $H_0$ at $\alpha = 0.05$
\item \textbf{Interpretation}: Complete rank separation --- on every tested equation, Hybrid v50\_2 extrapolation error is strictly less than Neural Network error ($U = 0$, $p = 1.11 \times 10^{-6}$)
\end{itemize}

\subsection{Effect Size Calculation (Cohen's d)}

\begin{align}
\mu_{\text{NN}} &= 1231.0\%, \quad \sigma_{\text{NN}} = 1842.0\% \quad \text{(CORRECTED)} \\
\mu_{\text{Hybrid}} &= 0.0\%, \quad \sigma_{\text{Hybrid}} = 0.0\% \\
s_{\text{pooled}} &= \sqrt{\frac{\sigma_{\text{NN}}^2 + \sigma_{\text{Hybrid}}^2}{2}} = 1302.5\% \\
d &= \frac{1231.0 - 0.0}{1302.5} = 0.95
\end{align}

\paragraph{Conservative Estimate:}
Using only NN standard deviation (most conservative approach):
\begin{equation}
d_{\text{conservative}} = \frac{1231.0}{1842.0} = 0.67 \quad \text{(medium-to-large effect)}
\end{equation}

\paragraph{Alternative Calculation (Glass's $\Delta$):}
Using only the control group (NN) standard deviation:
\begin{equation}
\Delta = \frac{1231.0 - 0.0}{1842.0} = 0.67
\end{equation}

\paragraph{Interpretation Guidelines:}
\begin{itemize}
\item $d = 0.2$: Small effect
\item $d = 0.5$: Medium effect
\item $d = 0.8$: Large effect
\item $d = 1.2$: Very large effect
\item Our result: $d = 0.67-0.95$ (medium-to-large effect with full rank separation on this benchmark)
\end{itemize}

\paragraph{Critical Note on Effect Size:} While the calculated Cohen's d appears moderate (0.67-0.95), this understates the true difference because:
\begin{enumerate}
\item Mann-Whitney $U = 0$ indicates full rank separation---every Hybrid error is strictly less than every NN error
\item Effect size calculations assume normal distributions, but NN errors are highly skewed (median = 86.7\%, mean = 1231\%)
\item The practical difference (near-zero vs 1,231\% mean error on this benchmark) is substantial and potentially critical for scientific applications
\end{enumerate}

\subsection{Interpolation Performance Comparison}

\paragraph{Kruskal-Wallis H Test:}
\begin{itemize}
\item \textbf{Test statistic}: $H = 42.3$
\item \textbf{P-value}: $p < 0.001$
\item \textbf{Interpretation}: Significant differences exist across methods
\end{itemize}

\paragraph{System Comparisons:}
\begin{table}[H]
\centering
\caption{Interpolation Performance Statistics by System}
\label{tab:interpolation_stats}
\begin{tabular}{lccccc}
\toprule
\textbf{System} & \textbf{n} & \textbf{Mean $R^2$} & \textbf{Median $R^2$} & \textbf{SD} & \textbf{Range} \\
\midrule
System 3 (LLM+Fallback) & 30 & 1.000 & 1.000 & 0.000 & [1.00, 1.00] \\
Hybrid v50\_2 & 15 & 0.931 & 1.000 & 0.147 & [0.58, 1.00] \\
Neural Network & 15 & 0.940 & 0.952 & 0.082 & [0.78, 1.00] \\
Pure PySR & 18 & 0.982 & 0.998 & 0.045 & [0.88, 1.00] \\
\bottomrule
\end{tabular}
\end{table}

\paragraph{Key Observation:} System 3 achieves R$^{2}$ = 1.000 on interpolation (n=30) but was not designed for extrapolation testing. Hybrid v50\_2's median (1.000) exceeds its mean (0.931) due to outlier presence, making median more representative of typical performance.

\subsection{Domain-Specific Success Rates}

Success rate defined as R$^{2}$ $> 0.95$:

\begin{table}[H]
\centering
\caption{Success Rate by Domain and Method}
\label{tab:domain_success_detailed}
\begin{tabular}{lccccc}
\toprule
\textbf{Method} & \textbf{Physics} & \textbf{Chemistry} & \textbf{Biology} & \textbf{DeFi} & \textbf{Economics} \\
\midrule
Hybrid v50\_2 & 100\% & 100\% & 67\% & 100\% & 67\% \\
Pure PySR & 100\% & 100\% & 100\% & 67\% & 67\% \\
LLM-Guided PySR & 100\% & 67\% & 67\% & 67\% & 67\% \\
Neural Network & 100\% & 67\% & 33\% & 67\% & 33\% \\
Pure LLM & 100\% & 33\% & 0\% & 33\% & 33\% \\
\midrule
\textbf{Avg across methods} & 100\% & 73\% & 53\% & 67\% & 53\% \\
\bottomrule
\end{tabular}
\end{table}

\paragraph{Correlation Analysis:} Pure LLM success rate strongly correlates with equation prevalence in web text (Spearman $\rho = 0.89$, $p < 0.001$)\citep{kambhampati2024llms}. This is consistent with the reproductive behaviour hypothesis: performance is higher where training-corpus exposure is greater.

\clearpage
\subsection{Validation Layer Effectiveness}

Error detection breakdown across four-layer validation framework:

\begin{table}[H]
\centering
\caption{Validation Layer Error Detection Statistics}
\label{tab:validation_stats}
\begin{tabular}{@{}lcccp{3cm}@{}}
\toprule
\textbf{Layer} & \textbf{Errors} & \textbf{\%} & \textbf{Cum.} & \textbf{Example} \\
\midrule
1. Dimensional Analysis & 18 & 12 & 12 & Energy = mass \\
2. Physics Constraints & 27 & 18 & 30 & Negative rate \\
3. Numerical Stability & 35 & 23 & 53 & Overflow \\
4. Prediction Accuracy & 70 & 47 & 100 & Wrong result \\
\midrule
\textbf{Total} & \textbf{150} & \textbf{100} & --- & --- \\
\bottomrule
\end{tabular}
\end{table}

\paragraph{Key Finding:} Single-layer validation (prediction accuracy only) would miss 53\% of errors caught by earlier layers on this benchmark, indicating that cascading validation contributed substantially to error detection for LLM-generated outputs.

\clearpage
% ============================================================================
% APPENDIX E: DISCOVERED EXPRESSION EXAMPLES
% ============================================================================


\section{Discovered Expression Examples}
\label{app:expressions}

\subsection{Perfect Recoveries}

\subsubsection{Kinetic Energy}

\begin{example}[Kinetic Energy Discovery]
\textbf{Ground truth}: $E = \frac{1}{2}mv^2$

\textbf{HypatiaX v50\_2 discovered}:
\begin{equation}
E = ((v \times m) \times v) \times 0.5 = 0.5mv^2
\end{equation}

\textbf{Analysis}: Exact recovery of ground truth. Discovered expression is algebraically identical.

\textbf{Interpolation}: $R^2 = 1.0000$ on training data (v $\in$ [1, 10] m/s)

\textbf{Extrapolation}: 0.0\% error at v = 50 m/s (5$\times$ training max)

\textbf{Physical interpretation}: Correctly captures quadratic dependence on velocity and linear dependence on mass.
\end{example}

\subsubsection{Ideal Gas Law}

\begin{example}[Ideal Gas Law Discovery]
\textbf{Ground truth}: $P = nRT/V$, where $R = 8.314$\,J/(mol$\cdot$K)

\textbf{HypatiaX v50\_2 discovered}:
\begin{equation}
P = n \times (T \times (8.314 / V))
\end{equation}

\textbf{Analysis}: Exact recovery including the universal gas constant $R = 8.314$ to full precision.

\textbf{Interpolation}: $R^2 = 1.0000$ on training data

\textbf{Extrapolation}: 0.0\% error across all pressure/temperature/volume ranges tested (up to 5$\times$ training range)

\textbf{Physical interpretation}: On this test case, symbolic regression recovered a constant consistent with the gravitational constant from data alone, without prior knowledge of its value\citep{lacava2021contemporary}.
\end{example}

\subsection{Near-Perfect Recoveries}

The Michaelis-Menten and Impermanent Loss equations were recovered with constant error $< 10^{-5}$, achieving $R^2 = 1.0000$ and near-zero extrapolation error (see the discovered-expression examples above).\textbf{[EDITORIAL NOTE: original text pointed to ``Listing~\texttt{\textbackslash ref\{sec:results\}}''; \texttt{sec:results} is the main paper's entire Results section, not a listing, so this appears to be a leftover/incorrect reference. Removed pending confirmation of intended target.]}

\subsection{Failure Case Analysis}

\subsubsection{Gravitational Force (Both Methods Failed)}

\begin{example}[Gravitational Force Failure]
\textbf{Ground truth}: $F = Gm_1m_2/r^2$, where $G = 6.674 \times 10^{-11}$\,N$\cdot$m$^2$/kg$^2$

\textbf{Neural Network Result}:
\begin{itemize}
\item \textbf{Interpolation}: $R^2 = 0.21$ (poor fit; relative error $> 100\%$, meeting the threshold in the main paper's \emph{Interpolation--Extrapolation Decoupling} definition)
\item \textbf{Extrapolation}: Not attempted due to interpolation failure
\item \textbf{Issue}: Gradient vanishing due to $10^{-11}$ scale
\item \textbf{Attempted solution}: Log-transform $\to$ marginal improvement ($R^2 = 0.35$)
\item \textbf{Root cause}: Optimization landscapes too flat at this scale for gradient descent
\end{itemize}

\textbf{HypatiaX v50\_2 Result}:
\begin{itemize}
\item \textbf{Interpolation}: $R^2 = -0.03$ (worse than baseline)
\item \textbf{Extrapolation}: Not attempted due to interpolation failure
\item \textbf{Discovered form}: $F = 1.23 \times 10^{13} \frac{m_1 m_2}{r^2}$ (correct functional form, wrong constant by $2 \times 10^{23}\times$)
\item \textbf{Issue}: PySR constant search range [$10^{-3}$, $10^3$] missed $G \sim 10^{-11}$
\item \textbf{Attempted solution}: Extend range to [$10^{-15}$, $10^{15}$] $\to$ no convergence in 50 iterations
\item \textbf{Root cause}: Search space too large for evolutionary algorithm to explore efficiently
\end{itemize}

\textbf{Recommended Solution}: 
\begin{enumerate}
\item \textbf{Dimensional analysis preprocessing}: Identify required constant scales before symbolic regression
\item \textbf{Logarithmic transformation}: Convert $F = Gm_1m_2/r^2$ to $\log F = \log G + \log m_1 + \log m_2 - 2\log r$, making the problem linear in log-space
\item \textbf{Adaptive constant ranges}: Dynamically adjust search based on data magnitude
\item \textbf{Multi-scale normalization}: Normalize inputs/outputs to [0.1, 10] range before discovery, then rescale
\end{enumerate}

\textbf{Lesson learned}: Very small constants ($< 10^{-6}$) require special handling in both neural and symbolic methods. This represents 1/15 test cases (6.7\% failure rate) and is a known limitation addressed in Section~\ref{suppb:limitations}.
\end{example}

\subsubsection{Henderson-Hasselbalch (Pure LLM Failure)}

\begin{example}[Henderson-Hasselbalch pH Calculation]
\textbf{Ground truth}: $\text{pH} = \text{pKa} + \log_{10}\left(\frac{[\text{A}^-]}{[\text{HA}]}\right)$

\textbf{Pure LLM Generated}:
\begin{equation}
\text{pH} = \text{pKa} \times \frac{[\text{A}^-]}{[\text{HA}]}
\end{equation}

\textbf{Failure Analysis}:
\begin{itemize}
\item \textbf{R$^{2}$ score}: $-12.3$ (predictions anticorrelated with ground truth)
\item \textbf{Errors}: 
  \begin{enumerate}
  \item Multiplication instead of addition
  \item Missing logarithm entirely
  \item Plausible but fundamentally wrong
  \end{enumerate}
\item \textbf{Root cause}: LLM training data likely contains more qualitative discussions of pH (''pH increases when conjugate base increases") than quantitative formulas
\item \textbf{Pattern matching failure}: LLM defaulted to simpler multiplicative relationship rather than logarithmic form
\end{itemize}

\textbf{Symbolic Rescue}:
PySR correctly discovered logarithmic relationship:
\begin{equation}
\text{pH} = \text{pKa} + \log_{10}\left(\frac{[\text{A}^-]}{[\text{HA}]}\right)
\end{equation}
Achieving $R^2 = 0.998$, demonstrating the value of hybrid architecture.

\textbf{Lesson learned}: LLM performance declined on this specialised equation, which is consistent with the main paper's \emph{LLM Reproductive Behaviour} finding---generation probability is theoretically predicted to be lower for expressions that are structurally distant from the training corpus\citep{kambhampati2024llms}.
\end{example}

\subsection{Complete Discovered Expressions Code Listing}

\begin{lstlisting}[caption=Discovered Symbolic Expressions (7 Perfect Recoveries), label=lst:discovered, language=Python]
# Kinetic Energy (PERFECT)
discovered: ((v * m) * v) * 0.5
ground_truth: 0.5 * m * v**2
R2: 1.0000, extrapolation_error: 0.0%

# Rate Law (PERFECT)
discovered: ((A_conc * 0.5) * A_conc) * B_conc  
ground_truth: 0.5 * A_conc**2 * B_conc
R2: 1.0000, extrapolation_error: 0.0%

# Ideal Gas Law (PERFECT)
discovered: n * (T * (8.314 / V))
ground_truth: n * 8.314 * T / V
R2: 1.0000, extrapolation_error: 0.0%

# Logistic Growth (PERFECT - algebraically equivalent)
discovered: ((N * -0.0003) + 0.3) * N
ground_truth: 0.3 * N * (1 - N/1000)
expanded: 0.3*N - 0.0003*N^2 (identical)
R2: 1.0000, extrapolation_error: 0.0%

# Price Impact (PERFECT)
discovered: swap / (swap + reserve)
ground_truth: dx / (x + dx)
R2: 1.0000, extrapolation_error: 0.0%

# Michaelis-Menten (NEAR-PERFECT)
discovered: (S / (S + 10.000007)) * 50.000008
ground_truth: 50 * S / (10 + S)
constant_error: <1e-5
R2: 1.0000, extrapolation_error: 0.0%

# Impermanent Loss (NEAR-PERFECT)
discovered: (2 * sqrt(r)) / (1 + r) - 1.000023
ground_truth: 2 * sqrt(r) / (1 + r) - 1
constant_error: 2.3e-5
R2: 1.0000, extrapolation_error: 0.0%
\end{lstlisting}

% ============================================================================
% APPENDIX F: REPRODUCIBILITY STATEMENT (CORRECTED 2026-03-21)
% ============================================================================


\section{Supplementary Materials}
\label{app:supplementary}

The following supplementary materials are available:

\begin{itemize}
\item \textbf{S1: Pure LLM Baseline Results} (\texttt{S1\_pure\_llm\_results.csv}) --- Complete results from Pure LLM baseline experiments showing 60\% success rate across 40 tests.

\item \textbf{S2: Extrapolation Test Data} (\texttt{S2\_extra\-polation\_data.csv}) --- Data comparing Hybrid v50\_2 extrapolation error (near-zero) versus Neural Network extrapolation error (1{,}231\% mean) on the Core 15 benchmark. Includes training R$^2$ scores, extrapolation errors at 1.2$\times$, 2$\times$, and 5$\times$ training ranges, and Mann-Whitney U-test results (U=0, p$<$10$^{-6}$).

\item \textbf{S3: Validation Logs} (\texttt{S3\_validation\_logs.json}) --- Complete validation reports for all 131 tests showing the four-layer validation framework in action. Each entry includes dimensional analysis, physics constraints, numerical stability checks, prediction accuracy metrics, timing data, and failure mode classification. Includes validation layer rejection rates (12\%, 18\%, 23\%, 47\%).

\item \textbf{S4: Discovered Symbolic Expressions} (\texttt{S4\_discovered\_expressions.\{txt,json\}}) --- All 131 discovered symbolic expressions in SymPy-compatible format, organized by domain (chemistry, biology, physics, DeFi, finance) with complete metadata including ground truth equations, performance metrics (R$^2$, timing, iterations), and algebraic equivalence verification.

\item \textbf{S5: Timing Breakdown} (\texttt{S5\_timing\_breakdown.csv}) --- Detailed computational cost analysis showing component-level timing: LLM initialization (0-39\%), symbolic search (50-91\%), and validation (7-11\%). Demonstrates the 4.3× speedup from LLM-guided initialization versus pure symbolic methods.

\item \textbf{S6: Domain-Specific Analysis} (\texttt{S6\_domain\_analysis.xlsx}) --- Statistical analysis by scientific domain showing success rates: Biology 100\%, Chemistry 100\%, Physics 100\%, DeFi 100\%, Finance 77.8\%. Includes per-equation breakdowns and domain-specific constraint validation results.

\item \textbf{S7: Failure Mode Taxonomy} (\texttt{S7\_failure\_mode\_taxonomy.pdf}) --- Comprehensive taxonomy of 23 distinct failure modes categorized by method (Pure LLM: 8 modes, Neural Network: 3 modes, Symbolic: 5 modes, Hybrid: 2 modes, Integration: 5 modes), severity (8 Critical, 4 High, 9 Medium, 2 Low), and detection layer. Includes detailed examples of silent semantic errors, large-scale extrapolation errors, distributional reasoning failures, and API misuse.

\item \textbf{S8: Figure Source Data} (\texttt{S8\_figures\_source\_data/}) --- Source data files for all figures in the paper:
\begin{itemize}
\item \texttt{figure1\_arrhenius\_extrapolation.csv} --- Arrhenius equation extrapolation comparison
\item \texttt{figure2\_domain\_comparison.csv} --- Success rates by scientific domain
\item \texttt{figure3\_validation\_breakdown.csv} --- Four-layer validation effectiveness
\item \texttt{figure4b\_domain\_breakdown.csv} --- Per-equation performance details
\item \texttt{figure5\_method\_comparison.csv} --- Overall method comparison
\item \texttt{figure\_5systems\_comparison.csv} --- Complete extrapolation data (1231\% vs 0\%)
\item \texttt{figure\_benchmark\_comparison.csv} --- Core 15 benchmark results
\item \texttt{figure\_timing\_comparison.csv} --- Speed-accuracy tradeoff analysis
\end{itemize}
\end{itemize}

All supplementary materials are provided in both synthetic versions (matching reported statistics exactly) and real experimental versions (direct from JSON logs) for maximum reproducibility and transparency. Complete documentation and usage examples are provided in accompanying README files.

\subsection{Supplementary Data Files}

All supplementary data are available at:
\begin{center}
{\small\url{https://github.com/sednabcn/LLM-HypatiaX-PAPERS}}
\end{center}

\begin{table}[H]
\centering
\caption{Supplementary Data Files}
\label{tab:supplementary_files}
\small
\begin{tabular}{@{}lp{8cm}@{}}
\toprule
\textbf{File} & \textbf{Contents} \\
\midrule
\texttt{S1\_pure\_llm\_results.csv} & Complete Pure LLM test results (n=40) with formulas, R$^{2}$ scores, failure modes \\
\texttt{S2\_extrapolation\_data.csv} & Extrapolation test data ($n=27$) at $1.2\times$, $2\times$, $5\times$ training ranges \\
\texttt{S3\_validation\_logs.json} & Full validation reports for all 131 tests \\
\texttt{S4\_discovered\_expressions.txt} & All discovered symbolic expressions in SymPy format \\
\texttt{S5\_timing\_breakdown.csv} & Detailed timing data by component and method \\
\texttt{S6\_domain\_analysis.xlsx} & Per-domain performance with statistical tests \\
\texttt{S7\_failure\_taxonomy.pdf} & Complete taxonomy of 23 failure modes with examples \\
\texttt{S8\_figures\_source\_data/} & Raw data for all 8 figures in CSV format \\
\bottomrule
\end{tabular}
\end{table}

\subsection{Reproducibility Materials}

Complete step-by-step tutorials for reproducing all experiments are available at:
\url{https://sednabcn.github.io/ai-llm-blog/tutorials/hypatiax/}

\begin{itemize}
\item \href{https://sednabcn.github.io/ai-llm-blog/tutorials/hypatiax/setup/}{Tutorial 1: Environment Setup and First Discovery} (15 min)
\item \href{https://sednabcn.github.io/ai-llm-blog/tutorials/hypatiax/experiments/}{Tutorial 2: Running Benchmark Experiments} (45 min active + 3--8 hours compute)
\item \href{https://sednabcn.github.io/ai-llm-blog/tutorials/hypatiax/analysis/}{Tutorial 3: Statistical Analysis and Publication Figures} (30 min)
\item \href{https://sednabcn.github.io/ai-llm-blog/tutorials/hypatiax/extensions/}{Tutorial 4: Custom Applications and Extensions} (45 min)
\end{itemize}

All source code and data generation scripts are available at:
\url{https://github.com/sednabcn/LLM-HypatiaX-PAPERS}

\subsection{Reproducibility Artifacts}

\paragraph{Docker Images:}
\begin{itemize}
\item \texttt{hypatiax:v1.0} - Full environment with all dependencies
\item \texttt{hypatiax:minimal} - Core symbolic regression only (no LLM)
\item \texttt{hypatiax:gpu} - GPU-accelerated version for neural network comparison
\end{itemize}

\paragraph{Source Code:}
All Python scripts and experimental protocols are available at:
\begin{center}
{\small\url{https://github.com/sednabcn/LLM-HypatiaX-PAPERS/tree/master/papers/2025-JMLR/hypatiax}}
\end{center}

 
\paragraph{Jupyter Notebooks:}
\begin{itemize}
\item \texttt{01\_data\_generation.ipynb} - Synthetic data generation for 15 equations
\item \texttt{02\_pure\_llm\_experiments.ipynb} - Pure LLM baseline (40 tests)
\item \texttt{03\_hybrid\_experiments.ipynb} - Hybrid v50\_2 experiments (30 tests)
\item \texttt{04\_extrapolation\_analysis.ipynb} - Extrapolation tests and statistical analysis
\item \texttt{05\_figure\_generation.ipynb} - Reproduction of all 8 figures
\item \texttt{06\_statistical\_tests.ipynb} - Mann-Whitney U, Kruskal-Wallis, effect sizes
\end{itemize}

All Python Jupyter notebooks are available at:
\begin{center}
{\footnotesize\url{https://github.com/sednabcn/LLM-HypatiaX-PAPERS/tree/master/papers/2025-JMLR/hypatiax/notebooks}}
\end{center}


\paragraph{Step-by-Step Tutorials:}
\begin{itemize}
\item \href{https://sednabcn.github.io/ai-llm-blog/tutorials/hypatiax/setup/}{Tutorial 1: Environment Setup and First Discovery} (15 min)
\item \href{https://sednabcn.github.io/ai-llm-blog/tutorials/hypatiax/experiments/}{Tutorial 2: Running Benchmark Experiments} (45 min active + 3--8 hours compute)
\item \href{https://sednabcn.github.io/ai-llm-blog/tutorials/hypatiax/analysis/}{Tutorial 3: Statistical Analysis and Publication Figures} (30 min)
\item \href{https://sednabcn.github.io/ai-llm-blog/tutorials/hypatiax/extensions/}{Tutorial 4: Custom Applications and Extensions} (45 min)
\end{itemize}
Available at: \url{https://sednabcn.github.io/ai-llm-blog/tutorials/hypatiax/}


\subsection{Additional Statistical Tests}

\begin{table}[H]
\centering
\caption{Additional Statistical Tests}
\label{tab:additional_stats}
{\small\begin{tabular}{@{}lccl@{}}
\toprule
\textbf{Comparison} & \textbf{Test} & \textbf{Result} & \textbf{Interpretation} \\
\midrule
Hybrid v50\_2 vs Pure PySR & Mann-Whitney & $U = 89$, $p = 0.18$ & No significant difference \\
System 3 vs Hybrid v50\_2 & Mann-Whitney & $U = 156$, $p = 0.03$ & System 3 better (interp.) \\
Classical vs DeFi (LLM) & Fisher's Exact & $p < 0.001$ & Domain effect significant \\
Run-to-run variability & Friedman & $\chi^2 = 34.2$, $p < 0.001$ & High variability \\
\bottomrule
\end{tabular}}
\end{table}

\paragraph{Power Analysis:}
\begin{itemize}
\item Extrapolation comparison (n=14 vs n=13): Power = 1.000 for d > 0.5
\item Domain comparisons (n=3 per domain): Power = 0.65 for d > 0.8
\item Method comparisons (n=15 per method): Power = 0.95 for d > 0.6
\end{itemize}

Recommendation: Future work should increase per-domain sample sizes to n $\geq$ 10 for robust inference on secondary effects.

% ============================================================================
% APPENDIX H: NOISE SENSITIVITY ANALYSIS
% ============================================================================


\section{Noise Sensitivity Analysis}
\label{app:noise_sensitivity}

To assess correspondence with Condition 3 of the main paper's \emph{LLM Reproductive Behaviour} finding (noise tolerance threshold $\sigma \leq 0.1 \cdot \text{std}(y)$), we conducted systematic experiments varying noise levels.

\subsection{Experimental Protocol}

We tested 15 classical formulas (see the main paper's Five-System Comparison table) with noise levels:
\[
\sigma \in \{0.01, 0.03, 0.05, 0.08, 0.10, 0.15, 0.20\} \times \text{std}(y)
\]

For each noise level, we:
\begin{enumerate}
\item Generated $n = 100$ training points with additive Gaussian noise: $y_{\text{noisy}} = y_{\text{true}} + \mathcal{N}(0, \sigma^2)$
\item Ran PySR with standard hyperparameters (Appendix~\ref{app:hyperparameters})
\item Evaluated success as $R^2 \geq 0.99$ and correct functional form
\item Repeated 5 times per configuration to account for evolutionary randomness
\end{enumerate}

\subsection{Results}

\begin{table}[H]
\centering
\caption{Success Rate vs Noise Level}
\label{tab:noise_sensitivity}
\begin{tabular}{@{}ccc@{}}
\toprule
\textbf{Noise Level ($\sigma / \text{std}(y)$)} & \textbf{Success Rate} & \textbf{95\% CI} \\
\midrule
0.01 & 100\% (75/75) & [95.2\%, 100\%] \\
0.03 & 97.3\% (73/75) & [90.7\%, 99.7\%] \\
0.05 & 93.3\% (70/75) & [85.1\%, 97.8\%] \\
0.08 & 78.7\% (59/75) & [67.7\%, 87.3\%] \\
0.10 & 66.7\% (50/75) & [54.8\%, 77.1\%] \\
0.15 & 45.3\% (34/75) & [33.8\%, 57.3\%] \\
0.20 & 24.0\% (18/75) & [14.9\%, 35.3\%] \\
\bottomrule
\end{tabular}
\end{table}

\subsection{Key Findings}

\paragraph{Sharp Threshold at $\sigma = 0.1$:}
Success rate drops from 93.3\% ($\sigma = 0.05$) to 45.3\% ($\sigma = 0.15$), consistent with the empirically derived threshold. Fisher's exact test shows a significant difference between $\sigma \leq 0.10$ and $\sigma > 0.10$ ($p < 0.001$).

\paragraph{Domain Variability:}
Classical mechanics formulas (kinetic energy, potential energy) tolerate higher noise ($\sigma \leq 0.15$) while chemistry formulas (Arrhenius, rate law) fail above $\sigma = 0.08$. This suggests domain-specific complexity affects noise sensitivity.

\paragraph{Failure Modes:}
Above threshold, common failures include:
\begin{itemize}
\item Incorrect exponents (e.g., $v^{1.8}$ instead of $v^2$ for kinetic energy)
\item Spurious additive constants to compensate for noise
\item Overfitting with complex expressions ($R^2 > 0.99$ on training, poor extrapolation)
\end{itemize}

\subsection{Implications}

Our main experiments use $\sigma = 0.05 \cdot \text{std}(y)$, safely below the threshold, ensuring 93.3\% baseline success rate from noise alone. Real-world applications should:
\begin{enumerate}
\item Estimate noise level from repeated measurements
\item Apply denoising if $\sigma > 0.08 \cdot \text{std}(y)$
\item Increase sample size by factor of $(\sigma / 0.05)^2$ to maintain success rate
\end{enumerate}

% ============================================================================
% APPENDIX I: EXTENDED DISCUSSION
% ============================================================================




\end{document}
